\section{Connection to the general FLRW visualization}
The earlier manuscript used a dimensionless scale factor $a(t_o)=1$
and a length-valued comoving radial coordinate $r$. Its radial metric is
\begin{equation}
 \dd s_{\mathrm{rad}}^2=-c^2\dd t^2+a^2(t)\dd r^2.
\end{equation}
For any positive plotting length $L$, define $R=La$ and $\psi=r/L$.
On an interval with $H>0$,
\begin{equation}\label{eq:plotting-metric}
 \dd s_{\mathrm{rad}}^2=-\frac{c^2}{H^2R^2}\dd R^2+R^2\dd\psi^2,
 \qquad \frac{\dd\psi}{\dd R}=-\frac{c}{HR^2}
\end{equation}
for an incoming future-directed ray. For a closed background with
present curvature radius $R_c$, setting $L=R_c$ gives
$R=A=R_ca$ and $\psi=\chi$. For flat or negatively curved spatial
geometry, the same polar map is a radial visualization, not a round
embedding of complete spatial slices. Changing arbitrary $L$ changes
angular sweep without changing the underlying physics.

For the illustrative history $a=(t/t_o)^p$ with $0<p<1$,
\begin{equation}
 \psi=\frac{ct_o}{L(1-p)}\left[1-(t/t_o)^{1-p}\right].
\end{equation}
The sweep is finite at the formal past limit. For $p=1$,
$\psi=(ct_o/L)\log(t_o/t)$ and the drawing is logarithmically spiral.
These shape statements are properties of the map and the chosen
history; neither implies physical rotation.

\section{Distances, horizons, and the all-sky early universe}
For a background comoving source,
\begin{equation}
 D_C=c\int_{t_e}^{t_o}\frac{\dd t}{a(t)},\qquad
 D_{\mathrm{proper}}(t)=a(t)D_C,\qquad
 1+z=\frac{1}{a(t_e)}.
\end{equation}
In a closed background, transverse comoving distance is
$D_M=R_c\sin(D_C/R_c)$ on the direct branch, and
$D_A=D_M/(1+z)$, $D_L=(1+z)D_M$ under the usual transparent-propagation
conditions \citep{Hogg1999}. These relations are background expressions;
the perturbed model requires optical transport to calculate their
direction-dependent corrections.

The particle horizon uses the convergent integral from the adopted past
boundary to reception. The event horizon uses the remaining future
integral when it converges. The Hubble radius $c/H$ is generally neither
of these causal horizons \citep{DavisLineweaver2004}.

In the thin-last-scattering background approximation, the emission
locus of the primary CMB photons received at $p_o$ is
\begin{equation}
 \mathcal S_*(p_o)=\partial J^-(p_o)\cap\Sigma_{t_*},
\end{equation}
restricted to the direct light cone before caustics. The causal past
$J^-$ itself also includes timelike-related events. The observed emission
locus samples the same early epoch in every direction; it is not a
material shell around the observer today. Finite last-scattering width,
scattering, and perturbations require a transfer calculation
\citep{SeljakZaldarriaga1996}. The geometry alone does not explain the
CMB spectrum, its near-uniformity, or the initial conditions.

\section{Numerical conventions and validity}
\begin{table}[htbp]
 \centering\small
 \begin{tabular}{@{}lp{0.72\linewidth}@{}}
 \toprule Symbol & Meaning\\\midrule
 $A$ & Length-valued background curvature radius\\
 $a$ & Dimensionless scale factor, used in the general-map appendix\\
 $\eta,s$ & Dimensionless conformal time and conformal lookback\\
 $\chi$ & Geodesic angular distance on the unit three-sphere\\
 $\alpha$ & Local observer-frame launch angle for a view line\\
 $\ph$ & Dimensionless metric potential $\Phi/c^2$\\
 $\Delta,\delta_N$ & Comoving and Newtonian-gauge density contrasts\\
 $V$ & Peculiar-velocity potential, $v^i/c=\nabla^i_\gamma V$\\
 $N,B$ & Lapse factor and conformal spatial factor of the truncated metric\\
 $\bm n,\bm p$ & Unit-sphere position and tangent optical momentum\\
 \bottomrule
 \end{tabular}
\end{table}
The default growth integration uses DOP853 with relative tolerance
$2\times10^{-12}$. Standard ray integration uses relative tolerance
$2\times10^{-10}$, checked against a tighter run. Source shooting uses
tangent-plane endpoint residuals and verifies the full embedded source
position. The numerical model checks its time interval; the example
also verifies small potential, density contrast, and peculiar velocity.
Changing width or amplitude requires rechecking all three, since a
small potential alone does not imply a linear matter perturbation.
